\documentclass[conference]{IEEEtran}
\\usepackage{fontspec}\\usepackage[english]{babel}\\babelprovide[import]{vietnamese}\\setmainfont{DejaVu Serif}
\usepackage{amsmath,amssymb,amsthm,mathtools,algorithm,algpseudocode,listings,xcolor,hyperref}
\usepackage{enumitem}\hypersetup{hidelinks}\lstset{language=Python,basicstyle=\footnotesize\ttfamily,breaklines=true,frame=single,columns=fullflexible}
\title{Minimax và Alpha--Beta Pruning: Phân tích chi tiết và ứng dụng vào Tic-Tac-Toe}
\author{\IEEEauthorblockN{Tên sinh viên}\IEEEauthorblockA{Bộ môn Trí tuệ nhân tạo\\Trường/Đơn vị}}
\begin{document}\maketitle
\begin{abstract}Bài báo trình bày Minimax và Alpha--Beta pruning cho trò chơi đối kháng hai người có tổng bằng không. Nội dung gồm mô hình trạng thái, hàm giá trị, công thức đệ quy, chứng minh trực giác về tính tối ưu, ý nghĩa của hai cận $\alpha,\beta$, độ phức tạp và ảnh hưởng của thứ tự nước đi. Tic-Tac-Toe 3x3 được giải đầy đủ: mô tả cây tìm kiếm, từng bước Selection--Expansion theo kiểu DFS, điều kiện cắt, và mã Python có thể chạy độc lập.\end{abstract}
\begin{IEEEkeywords}Minimax, alpha-beta pruning, game tree, Tic-Tac-Toe, adversarial search\end{IEEEkeywords}
\section{Giới thiệu}Trong trò chơi đối kháng, một nước tốt phải xét phản hồi tốt nhất của đối thủ. Minimax duyệt cây trò chơi đến trạng thái kết thúc rồi truyền giá trị ngược lên: MAX chọn lớn nhất, MIN chọn nhỏ nhất. Alpha--Beta giữ nguyên kết quả Minimax nhưng bỏ qua các nhánh không thể thay đổi quyết định. Với Tic-Tac-Toe, không cần heuristic hay giới hạn độ sâu vì cây hữu hạn và đủ nhỏ để giải đến terminal.
\section{Mô hình trò chơi và giả thiết}Trạng thái là $s=(b,p)$, trong đó $b\in\{X,O,\cdot\}^9$ và $p$ là người sắp đi. $A(s)$ là các ô trống; $T(s,a)$ đặt quân $p$ vào ô $a$ và đổi lượt. $Z$ gồm trạng thái có ba quân liên tiếp hoặc bàn đầy. Giả thiết chuẩn là thông tin hoàn hảo, hai bên đi luân phiên, không ngẫu nhiên, và tổng phần thưởng bằng không. Dùng tiện ích từ góc nhìn người chơi gốc: thắng $+1$, hòa $0$, thua $-1$.
\section{Minimax}Gọi $V(s)$ là giá trị tối ưu từ trạng thái $s$. Khi $s$ terminal, $V(s)=U(s)$. Nếu lượt của MAX, $V(s)=\max_{a\in A(s)}V(T(s,a))$; nếu lượt của MIN, $V(s)=\min_{a\in A(s)}V(T(s,a))$. Nước tối ưu là $a^*=\arg\max_a V(T(s,a))$ tại nút gốc của MAX.\begin{algorithm}[h]\caption{Minimax đệ quy}\begin{algorithmic}[1]\Require trạng thái $s$, người chơi gốc $r$\If{$s$ kết thúc}\State \Return $U(s,r)$\EndIf\If{$player(s)=r$}\State \Return $\max_{a\in A(s)}\mathrm{Minimax}(T(s,a),r)$\Else\State \Return $\min_{a\in A(s)}\mathrm{Minimax}(T(s,a),r)$\EndIf\end{algorithmic}\end{algorithm}
\subsection{Vì sao đúng}Tại terminal, giá trị là định nghĩa kết quả. Giả sử các nút con đã có giá trị đúng. MAX có thể chọn đúng một con, nên chọn giá trị lớn nhất; MIN ép MAX nhận giá trị nhỏ nhất. Quy nạp theo độ sâu chứng minh giá trị tại gốc là kết quả khi cả hai chơi tối ưu. Nếu bàn không còn nước thắng, giá trị gốc bằng $0$: máy không thua, nhưng cũng không thể ép thắng.
\section{Alpha--Beta pruning}Trong DFS, $\alpha$ là giá trị tốt nhất MAX đã đảm bảo trên đường hiện tại; $\beta$ là giá trị tốt nhất MIN đã đảm bảo. Tại nút MAX, cập nhật $\alpha\leftarrow\max(\alpha,v)$; tại MIN, cập nhật $\beta\leftarrow\min(\beta,v)$. Khi $\alpha\ge\beta$, dừng vòng lặp: một tổ tiên đã có lựa chọn tốt hơn hoặc bằng, vì vậy phần còn lại không thể được chọn.\begin{algorithm}[h]\caption{Alpha--Beta}\begin{algorithmic}[1]\Require $s,\alpha=-\infty,\beta=+\infty$\If{$s$ kết thúc}\State \Return $U(s)$\EndIf\If{MAX}\State $v\gets-\infty$\For{$a\in A(s)$}\State $v\gets\max(v,AB(T(s,a),\alpha,\beta))$; $\alpha\gets\max(\alpha,v)$\If{$\alpha\ge\beta$}\State \textbf{break}\EndIf\EndFor\Else\State $v\gets+\infty$\For{$a\in A(s)$}\State $v\gets\min(v,AB(T(s,a),\alpha,\beta))$; $\beta\gets\min(\beta,v)$\If{$\alpha\ge\beta$}\State \textbf{break}\EndIf\EndFor\EndIf\State \Return $v$\end{algorithmic}\end{algorithm}
\paragraph{Bảo toàn kết quả.} Cắt chỉ xảy ra sau khi đã có một cận đủ mạnh. Các giá trị chưa duyệt chỉ có thể làm nút hiện tại tốt hơn cho người đang tối ưu, nhưng điều đó đã bị chặn bởi $\beta$ hoặc $\alpha$ ở tổ tiên; vì vậy chúng không đổi giá trị minimax trả về.
\section{Phân tích độ phức tạp}Với branching factor $b$ và độ sâu $d$, Minimax cần $O(b^d)$ nút và bộ nhớ DFS $O(d)$ (ngoài biểu diễn trạng thái). Alpha--Beta có cùng cận xấu nhất, nhưng nếu thứ tự nước đi gần tối ưu, số nút có thể giảm xuống $O(b^{d/2})$. Vì vậy move ordering quyết định hiệu quả. Trong Tic-Tac-Toe, duyệt ô trung tâm, góc, rồi cạnh thường giúp cắt sớm; mã minh họa duyệt theo chỉ số để minh bạch.
\section{Cơ chế search trên Tic-Tac-Toe}Mỗi lần gọi hàm, thuật toán thực hiện: (1) kiểm tra thắng/hòa; (2) sinh toàn bộ ô trống; (3) đặt quân và đệ quy đổi lượt; (4) nhận giá trị từ lá; (5) cập nhật max hoặc min; (6) với Alpha--Beta, cập nhật cận và có thể cắt. Ví dụ bàn rỗng, X thử ô 1, sau đó O thử các ô còn lại; mỗi nhánh được mở rộng cho đến khi có dòng thắng hoặc đủ 9 quân. Khi một nhánh cho O giá trị $-1$, tại nút O phép MIN giữ $-1$ và truyền lên; nếu tổ tiên MAX đã có $\alpha=0$, mọi nút MIN có $\beta\le0$ sẽ cắt phần còn lại khi $\alpha\ge\beta$. Kết quả tối ưu của bàn rỗng là hòa ($0$); một máy tối ưu không thể bị đánh bại.
\subsection{Bản ghi search có thể kiểm chứng}Mã đếm `nodes` cho số trạng thái đã thăm và `cutoffs` cho số lần cắt. Với mỗi nước ứng viên ở gốc, chương trình in giá trị $-1,0,+1$. `--show-search` in giá trị từng ô để quan sát các bước đánh giá. Phần terminal kiểm tra cả tám đường thắng; không dùng độ sâu cố định nên kết quả là tối ưu tuyệt đối.
\section{Mã Python}\lstinputlisting{ttt_minimax.py}
\section{Kết luận}Minimax cung cấp quyết định tối ưu bằng cách mô hình hóa phản hồi đối kháng. Alpha--Beta cho cùng đáp án nhưng giảm công việc nhờ cận và thứ tự duyệt. Tic-Tac-Toe là ví dụ lý tưởng để kiểm chứng vì có thể tìm toàn bộ cây; với trò chơi lớn hơn cần thêm giới hạn độ sâu, hàm heuristic, bảng chuyển trạng thái và move ordering.
\begin{thebibliography}{9}\bibitem{russell}S. Russell and P. Norvig, \emph{Artificial Intelligence: A Modern Approach}, 4th ed., Pearson, 2021.\bibitem{knuth}D. E. Knuth and R. W. Moore, ``An analysis of alpha-beta pruning,'' \emph{Artificial Intelligence}, vol. 6, pp. 293--326, 1975.\bibitem{shannon}C. Shannon, ``Programming a computer for playing chess,'' \emph{Philosophical Magazine}, vol. 41, pp. 256--275, 1950.\end{thebibliography}\end{document}
